\begin{abstract}
Training-free methods let a diffusion transformer trained near one megapixel generate $4096\times4096$ images by editing its rotary position embedding (RoPE). Methods are often compared as complete packages, leaving the contribution of each edit unclear. On \krea, over \N{ho@nprompts} held-out prompts with paired bootstrap intervals, \sega's base and global temperature alone recover all of its ImageReward gain over an NTK base (\shp{ir}, interval \shci{ir}). They tie the released method on all three reward models and crop MUSIQ, where ``tie'' means that the difference interval covers zero, not equivalence. The per-band allocation adds \hoa{alloc}{musiq} full-frame MUSIQ points and reduces excess fine-scale power. Four per-band laws that we derive from efficient coding and contrast constancy do no better: the two evaluated held out tie or lose to \sega{} on every reward and image-quality metric, and trajectory diagnostics suggest why. A development comparison also finds a halftone lattice on \fluxdev{} under official \dype{} code but not on \krea{}; differing guidance prevents attributing this contrast to the checkpoint alone. \sega's printed formulation and released code differ at each of the six stages we compared. Finally, 224-pixel reward models cannot directly resolve the top \N{nyq@nbands} octaves of a 4K image, and the default flow-shift extrapolation collapses the noise schedule. Claims for spectral or content-adaptive RoPE edits should be tested against a base-plus-temperature control with full-resolution metrics.
\end{abstract}
